% MMW-METHOD-CONTRACT: ee95080ed721168d72ddba0989ae7188b6e8baaceecbefd9c5b20fd61ea1aadf
\section{模型的建立与求解}

\subsection{数据预处理与统一计量口径}

首先，工作负载、GPU 中心、网络时延、功率映射、区域逐时数据与储能信息被按字段说明进行一致性核验。工作负载共包含 $50000$ 条记录，全部通过输入校验，缺失单元格数量为 $0$。区域、任务类型和小时索引被统一编码，任务类型与单位等效 GPU 功率严格按照功率映射表连接；区域间通信采用有向单向时延矩阵，不利用对称性补齐或替换原始记录。逐时电力数据覆盖任务主时域、末端结清时域和终端结算时点，其中第 $2406$ 小时被禁止安排任务，只用于能源与储能状态结算。

由于任务持续时间以分钟计量，而资源约束和能源参数以小时给出，任务 $i$ 的持续小时数及有效截止时点被定义为
\begin{equation}
h_i=\frac{d_i}{60},\qquad
F_i=\min\{f_i^{\mathrm{valid}},2406\}.
\label{eq:duration_deadline}
\end{equation}
若任务在整数时刻 $s$ 开工，则其与小时区间 $[t,t+1)$ 的实际重叠时长及瞬时占用指示量分别为
\begin{align}
\omega_{it}(s)
&=\max\{0,\min(t+1,s+h_i)-\max(t,s)\}, \label{eq:overlap}\\
\chi_{it}(s)
&=\mathbf 1\{\omega_{it}(s)>0\}. \label{eq:occupancy}
\end{align}
式~\eqref{eq:overlap} 被用于 GPU-hour 和小时电量计算，式~\eqref{eq:occupancy} 被用于 GPU、IT 功率与设施功率的瞬时容量审计。由此，非整小时任务不会被错误地按整小时计费，也不会因为采用小时平均值而掩盖短时容量冲突。

系统成本、购电碳排放、有效网络时延和新能源利用率统一按
\begin{align}
C_{\mathrm{op}}
&=\sum_{r,t}\left(c_{rt}^{buy}G_{rt}^{buy}
-c_{rt}^{sell}G_{rt}^{sell}\right)\Delta t, \label{eq:cost}\\
E_{\mathrm{CO2}}
&=\sum_{r,t}\kappa_{rt}G_{rt}^{buy}\Delta t, \label{eq:carbon}\\
L_{\mathrm{net}}
&=\frac{\sum_i g_ih_i\ell_{o_ir_i}}{\sum_i g_ih_i}, \label{eq:latency}\\
U_{\mathrm{RE}}
&=\frac{\sum_{r,t}(R_{rt}^{use}+R_{rt}^{ch}+R_{rt}^{sell})\Delta t}
{\sum_{r,t}R_{rt}^{av}\Delta t} \label{eq:renewable_utilization}
\end{align}
复算。式~\eqref{eq:latency} 和式~\eqref{eq:renewable_utilization} 仅在各自分母为正时计算；否则对应指标被标记为不可用，并退出约束、评分和排序，禁止使用任意小常数替代分母。多目标候选仅在通过全部硬约束后才参与比较，其指标差异采用有限正中位绝对尺度标准化。该处理可缓解成本、碳排放、时延和新能源利用率量纲不同造成的评分偏置，方法思想与稳健统计及多目标决策的基本原则一致\cite{huber1964robust,miettinen1999nonlinear}。

% MMW-ALGORITHM: 固定Huber滞后回归；50000任务NumPy区间切片确定性贪心调度；按真实调度重建AI小时电量和设施负荷；问题二从真实附件中定位区域级MaxGridImport_MW和MaxGridExport_MW并执行购售电边界及互斥审计；问题三使用Baseline AI负荷执行无储能比较与规则型SOC前向策略；问题四使用调度派生FacilityLoad执行联合储能前推并复算情景；MMW_PILOT只读取真实输入并执行有限样本门禁

实际程序采用 Python 与 scikit-learn 实现。需求预测采用固定 Huber 滞后回归，任务调度采用 NumPy 区间切片支持的确定性贪心可行构造，能源流由真实调度形成的小时电量重新计算。问题二执行新能源优先分流和购售电边界审计，问题三采用基准 AI 负荷下的规则型 SOC 前向策略，问题四依据调度派生设施负荷重新前推储能状态并复算各情景。除新能源波动情景使用固定随机种子 $2026$ 外，其余构造均为确定性过程。

\subsection{问题一：需求预测与基础算力调度}

\subsubsection{GPU 需求预测模型}

% MMW-ID: OBJ-Q1
区域 $r$、任务类型 $k$ 在小时 $t$ 的到达 GPU 需求被聚合为
\begin{equation}
D_{rkt}=\sum_{i:o_i=r,\;k_i=k,\;a_i=t}g_i.
\label{eq:demand_aggregation}
\end{equation}
预测模型同时引入短期、日内与周内滞后项，并采用
\begin{equation}
z_t=\left[
\sin\frac{2\pi t}{24},
\cos\frac{2\pi t}{24},
\sin\frac{2\pi t}{168},
\cos\frac{2\pi t}{168}
\right]^{\mathsf T}
\label{eq:periodic_features}
\end{equation}
表示周期特征。未截断的回归输出为
\begin{equation}
u_{rkt}=\beta_0+
\sum_{\ell\in\{1,2,3,24,48,168\}}
\beta_\ell D'_{rk,t-\ell}
+\gamma^{\mathsf T}z_t+\alpha_r+\delta_k.
\label{eq:huber_predictor}
\end{equation}
Huber 损失兼具平方损失在小残差区间内的平滑性和绝对损失对极端残差的稳健性\cite{huber1964robust}。实际程序将超参数固定为 $\epsilon=1.35$、$\alpha=0.001$，并在模型拟合完成后执行
\begin{equation}
\widehat D_{rkt}=\max\{0,u_{rkt}\},
\label{eq:nonnegative_prediction}
\end{equation}
从而保证预测需求非负，而不改变训练阶段的凸损失结构。

% MMW-ID: CON-Q1-6
预测误差采用
\begin{align}
\mathrm{MAE}&=\frac{1}{n}\sum_j|y_j-\widehat y_j|,\\
\mathrm{RMSE}&=\sqrt{\frac{1}{n}\sum_j(y_j-\widehat y_j)^2}
\end{align}
进行评价。WAPE 仅在真实值绝对和为正时计算，否则返回结构化不可用状态。最终测试结果表明，区域—任务类型序列的总体 MAE 为 $24.973340521042196$ GPU，总体 RMSE 为 $35.69449198489753$ GPU。测试区间共形成 $432$ 条区域—任务类型逐小时预测记录，需求统计明细共形成 $43200$ 行结果。

% MMW-LIMITATION: D2
方法契约原计划比较季节性基线并执行完整超参数网格，但实际运行采用固定 Huber 超参数，未完成全部候选的验证选模。因此，上述误差反映已实现固定模型在保留测试区间上的预测能力，而不能被解释为候选模型集合中的最佳验证结果。

\subsubsection{任务合法候选域与唯一执行约束}

% MMW-ID: CON-Q1-1
任务 $i$ 的合法执行区域—开工时刻集合被定义为
\begin{equation}
\Omega_i=\left\{(r,s):
s\ge a_i,\;
s+h_i\le F_i,\;
\ell_{o_ir}\le\ell_i^{\max},\;
\omega_{i,2406}(s)=0
\right\}.
\label{eq:candidate_domain}
\end{equation}
若 $\Omega_i$ 为空，则任务及其逐项排除原因应被记录，当前构造被结构化终止，不得通过放宽时延、截止时间或终点边界强行生成方案。

% MMW-ID: CON-Q1-2
每个任务必须且只能选择一个候选，实时推理任务的开工时刻必须等于到达时刻，即
\begin{align}
\sum_{(r,s)\in\Omega_i}x_{irs}&=1,\\
s_i&=\sum_{(r,s)\in\Omega_i}sx_{irs},\qquad c_i=s_i+h_i,\\
s_i&=a_i,\qquad i\in\mathcal I_{\mathrm{RT}}.
\end{align}
该约束同时保证任务不被拆分或迁移，并通过候选域保证任务不早于到达、不晚于截止且不占用第 $2406$ 小时。

\subsubsection{算力和功率容量约束}

% MMW-ID: CON-Q1-3
区域逐时 GPU、IT 功率和设施功率分别满足
\begin{align}
\sum_{i,s}g_i\chi_{it}(s)x_{irs}
&\le G_r^{\max}, \label{eq:gpu_capacity}\\
L_{rt}^{N}+\sum_{i,s}g_ip_{k_i}^{GPU}\chi_{it}(s)x_{irs}
&\le P_r^{IT,\max}, \label{eq:it_capacity}\\
\pi_r\left(
L_{rt}^{N}+\sum_{i,s}g_ip_{k_i}^{GPU}\chi_{it}(s)x_{irs}
\right)&\le P_r^{F,\max}. \label{eq:facility_capacity}
\end{align}
非 AI IT 负荷在式~\eqref{eq:it_capacity} 与式~\eqref{eq:facility_capacity} 中均被保留，因而 PUE 作用于区域总 IT 功率，而不是只作用于可迁移的 AI 部分。

% MMW-ID: CON-Q1-4
能源结算改用实际重叠时长：
\begin{align}
E_{rt}^{AI}
&=\sum_{i,s}g_ip_{k_i}^{GPU}\omega_{it}(s)x_{irs},\\
L_{rt}^{AI,avg}
&=\frac{E_{rt}^{AI}}{\Delta t},\\
L_{rt}^{F,avg}
&=\pi_r(L_{rt}^{N}+L_{rt}^{AI,avg}).
\end{align}
因此，容量审计使用 $\chi_{it}(s)$，能源核算使用 $\omega_{it}(s)$，两类口径在同一任务方案上被分别计算。

% MMW-ID: CON-Q1-5
区域逐时 GPU 利用率同时给出峰值口径和 GPU-hour 平均口径：
\begin{align}
U_{rt}^{peak}
&=\frac{\sum_{i,s}g_i\chi_{it}(s)x_{irs}}{G_r^{\max}},\\
U_{rt}^{avg}
&=\frac{\sum_{i,s}g_i\omega_{it}(s)x_{irs}}
{G_r^{\max}\Delta t}.
\end{align}
只有 $G_r^{\max}>0$ 时才执行除法。甘特图和容量审计采用 $U_{rt}^{peak}$，负荷描述采用 $U_{rt}^{avg}$。最后 $24$ 小时的 GPU 利用情况见图~\ref{fig:gpu_utilization}。

\begin{figure}[htbp]
\centering
\includegraphics[width=0.88\textwidth]{fig_1_gpu_utilization.png}
\caption{最后 $24$ 小时各区域 GPU 利用率}
\label{fig:gpu_utilization}
\end{figure}

图~\ref{fig:gpu_utilization} 按统一时间轴呈现了末端调度中的区域资源占用变化，可用于识别任务集中到达所形成的短时负荷波动。图中利用率应与任务甘特区间共同解释，因为同一小时内的峰值占用和按重叠时长计算的平均占用具有不同含义。容量是否合规最终以逐任务区间重叠和瞬时 GPU 审计为准，不能仅凭曲线的小时平均形态作出判断。

\subsubsection{实际调度算法与结果}

任务先按实时属性、有效截止、持续时间、GPU 需求和任务编号构成固定词典序，再在完整候选域内依据等待时间、单向时延、加入任务后的峰值 GPU 利用率、区域及开工时刻依次尝试。每次放置后，式~\eqref{eq:gpu_capacity}--\eqref{eq:facility_capacity} 均被重新审计。最后 $24$ 小时共有 $538$ 个实际到达任务获得执行区域及开工、完成时刻，其中 $68$ 个跨越第 $2399$ 小时的任务在第 $2400$ 至第 $2405$ 小时完成结清。

% MMW-ID: CON-Q1-7
理论 formulation 为候选耗尽时在局部冲突集内进行有限回溯，并在达到 $N_{\mathrm{back}}=2000$ 或搜索截止仍未形成完整方案时返回 \texttt{feasibility\_unresolved}。实际实现未执行该回溯，而是保留确定性贪心可行构造；因此，若复杂实例中贪心过程失败，只能说明已声明构造未获得完整方案。

% MMW-LIMITATION: D1
实际任务调度没有实现有限冲突回溯，故其探索范围小于数学模型规定的完整候选回溯范围。

% MMW-LIMITATION: D8
正式任务构造采用内部时间截止，为情景复算、审计和文件输出保留运行时间。达到该截止而尚未完成时，应返回 \texttt{feasibility\_unresolved}，不得将部分任务方案作为正式结果。

% MMW-LIMITATION: D10
试跑阶段仅对真实输入和有限任务样本执行合法时间窗门禁；试跑通过只说明输入接口与局部候选构造可运行，不代表全时域调度已经由试跑证明可行。

\subsection{问题二：无储能碳感知任务调度}

\subsubsection{任务负荷和能源分流模型}

% MMW-ID: OBJ-Q2
问题二的 formulation 从问题一可行方案出发，以单任务迁移或开工时刻变更构造候选。候选的成本、碳排放、有效网络时延、新能源未利用率和等待量应采用相对基准真实差值的有限正中位数尺度标准化，且只在完整轨迹通过硬约束后参与评分。

% MMW-ID: CON-Q2-1
任务侧继续执行唯一选址、实时即时开工、到达与截止、单向时延、瞬时 GPU、IT 功率、设施功率和真实重叠电量约束。问题二以全部 $50000$ 个实际任务重新形成调度负荷，而非采用问题一的预测需求替代真实任务。

% MMW-ID: CON-Q2-2
问题二不配置储能，故
\begin{equation}
R_{rt}^{ch}=G_{rt}^{ch}=C_{rt}=D_{rt}=0.
\end{equation}
新能源守恒与负荷平衡满足
\begin{align}
R_{rt}^{use}+R_{rt}^{sell}+R_{rt}^{curt}
&=R_{rt}^{av}, \label{eq:q2_renewable}\\
G_{rt}^{buy}+R_{rt}^{use}
&=L_{rt}^{F,avg}, \label{eq:q2_balance}\\
G_{rt}^{sell}&=R_{rt}^{sell}. \label{eq:q2_sell}
\end{align}

% MMW-ID: CON-Q2-3
购售电边界按真实附件中的区域级进口、出口上限确定：
\begin{align}
0\le G_{rt}^{buy}
&\le I_r^{\max}y_{rt}^{grid},\\
0\le G_{rt}^{sell}
&\le X_r^{grid}(1-y_{rt}^{grid}),
\qquad y_{rt}^{grid}\in\{0,1\}.
\end{align}
问题二不读取储能 SOC，也不使用储能附件中的售电限制，从而避免将问题三、四的储能边界错误施加于无储能调度。

% MMW-ID: CON-Q2-4
实际能源流采用确定性新能源优先规则：
\begin{align}
R_{rt}^{use}&=\min\{R_{rt}^{av},L_{rt}^{F,avg}\},\\
G_{rt}^{buy}&=L_{rt}^{F,avg}-R_{rt}^{use}.
\end{align}
当新能源不足时，区域不得同时售电；当新能源富余时，外送量受 $X_r^{grid}$ 限制，其余部分计为弃电。若负荷缺口超过区域购电上限，对应任务候选即被判为不可行。

% MMW-ID: CON-Q2-5
成本和碳排放分别由式~\eqref{eq:cost} 与式~\eqref{eq:carbon} 重新计算。结果显示，问题二运行成本为 $-419658147.7075976$ 元，碳排放为 $14.391601534493985$ tCO$_2$。负成本表示按既定购售电核算口径得到的售电收入超过购电支出，不能将其直接解释为系统不存在运行费用。

% MMW-ID: CON-Q2-6
网络时延和新能源利用率分别执行正分母门禁。有效网络时延为 $5.159384110987204$ ms，新能源利用率为 $0.6790780794609942$。任务方案中有 $209$ 个任务的执行区域不同于来源区域，第 $2400$ 至第 $2405$ 小时完成结清的末端弹性任务为 $68$ 个。GPU 容量、IT 功率、设施功率、网络时延和完成时限的统一验证结果通过。

% MMW-LIMITATION: D3
实际程序未执行 formulation 中规定的单任务局部改进，仅对确定性任务方案进行了完整能源复算与约束审计。因此，$209$ 个迁移任务来源于已实现贪心方案，而不是局部搜索收敛后的迁移结果。

\subsection{问题三：固定负荷下的储能协同模型}

\subsubsection{固定设施负荷与能源平衡}

% MMW-ID: OBJ-Q3
问题三固定采用基准 AI IT 负荷和非 AI IT 负荷。数学目标拟在成本、碳排放、非负峰值净购电、绝对爬坡和有效新能源未利用率之间进行标准化权衡；无储能状态只作为比较对象，不作为含储能方案能否启动的判据。

% MMW-ID: CON-Q3-1
区域设施负荷和能源流满足
\begin{align}
L_{rt}^{F,avg}
&=\pi_r(L_{rt}^{BAI}+L_{rt}^{N}), \label{eq:q3_load}\\
R_{rt}^{use}+R_{rt}^{ch}+R_{rt}^{sell}+R_{rt}^{curt}
&=R_{rt}^{av}, \label{eq:q3_re_balance}\\
G_{rt}^{load}+R_{rt}^{use}+D_{rt}
&=L_{rt}^{F,avg}, \label{eq:q3_load_balance}\\
C_{rt}&=R_{rt}^{ch}+G_{rt}^{ch},\\
G_{rt}^{buy}&=G_{rt}^{load}+G_{rt}^{ch},\\
G_{rt}^{sell}&=R_{rt}^{sell}.
\end{align}
新能源被明确拆分为直接消纳、储能充电、外送和弃电，从而避免将外送误认为新能源平衡残差。

\subsubsection{SOC 动态和运行边界}

% MMW-ID: CON-Q3-2
初始 SOC 是第 $0$ 小时运行前的唯一状态，储能递推为
\begin{align}
S_{r,-1}&=S_r^0,\\
S_{rt}&=S_{r,t-1}
+\eta_r^cC_{rt}\Delta t
-\frac{D_{rt}\Delta t}{\eta_r^d}, \label{eq:soc}\\
S_r^{min}&\le S_{rt}\le S_r^{max},\\
S_{r,2406}&\ge S_r^0. \label{eq:terminal_soc}
\end{align}
该定义保持 SOC 的绝对量纲，不把第 $0$ 小时重新归零，也不使用逐时基准 SOC 替代优化状态。

% MMW-ID: CON-Q3-3
充放电功率和购售电行为满足
\begin{align}
0\le C_{rt}&\le C_r^{\max},&
0\le D_{rt}&\le D_r^{\max},\\
0\le G_{rt}^{buy}
&\le I_r^{\max}y_{rt}^{grid},&
0\le G_{rt}^{sell}
&\le \min\{X_r^{grid},X_r^{storage}\}(1-y_{rt}^{grid}).
\end{align}
每小时储能动作只允许空闲、充电或放电之一，购电和售电亦保持互斥，以排除虚假能源循环和无约束套利。

% MMW-ID: CON-Q3-4
无储能状态按 $C_{rt}=D_{rt}=0$ 构造并用于比较。若其违反进口上限，应记录 \texttt{no\_storage\_baseline\_infeasible}，但含储能方案仍须从 $S_{r,-1}=S_r^0$ 独立前推。只有含储能构造未取得完整可行轨迹时，才返回 \texttt{feasibility\_unresolved}。

% MMW-ID: CON-Q3-5
每次储能动作后均检查终端 SOC 的必要可达性：
\begin{equation}
S_{rt}+\sum_{u=t+1}^{2406}
\eta_r^cC_r^{\max}\Delta t\ge S_r^0.
\end{equation}
随后，购电上限、能源守恒、SOC 边界、充放电互斥、购售电互斥及终端状态被按完整时域重新验证。

% MMW-ID: CON-Q3-6
净购电、峰值与绝对爬坡量定义为
\begin{align}
N_{rt}&=G_{rt}^{buy}-G_{rt}^{sell},\\
P_r^{peak}&=\max\{0,\max_tN_{rt}\},\\
V_r&=\sum_{t=1}^{2406}|N_{rt}-N_{r,t-1}|.
\end{align}
等效完整循环仅在储能容量为正时计算，否则返回不适用状态。

% MMW-ID: CON-Q3-7
新能源利用率继续采用式~\eqref{eq:renewable_utilization} 的正分母门禁，分母不正时不进入储能候选评分。

\subsubsection{规则型前向求解结果}

实际程序先计算无储能比较状态，再采用逐时规则型 SOC 前向策略构造含储能方案。无储能运行成本为 $-419820327.7393099$ 元，含储能运行成本为 $-419516103.81021965$ 元，变化量为 $304223.92909026146$ 元。这说明在已实现规则下，储能并未带来成本下降，因而不能把该结果表述为经济性改善。

无储能和含储能碳排放均为 $8.850562571909007$ tCO$_2$，变化量为 $0.0$ tCO$_2$；六区域非负峰值净购电功率之和均为 $31.597867090000022$ MW，变化量为 $0.0$ MW；绝对爬坡量均为 $63.195734180000045$ MW，变化量为 $0.0$ MW。含储能方案累计充电量为 $2076.3154884465794$ MWh，累计放电量为 $0.0$ MWh，第 $2406$ 小时终端 SOC 合计为 $3540.0$ MWh。相应 SOC 轨迹如图~\ref{fig:storage_soc} 所示。

\begin{figure}[htbp]
\centering
\includegraphics[width=0.88\textwidth]{fig_2_storage_soc.png}
\caption{各区域储能 SOC 轨迹}
\label{fig:storage_soc}
\end{figure}

图~\ref{fig:storage_soc} 展示了规则型前向策略下储能状态随时间的变化过程，轨迹必须结合初始 SOC、上下限和终端约束共同判断。累计放电量为 $0.0$ MWh，说明本次规则型方案没有形成依靠放电削峰或套利的循环行为。由于累计充电量为正而成本相对无储能状态增加，所得轨迹更接近终端能量恢复或新能源吸纳规则的结果，不能作为储能经济调度的一般性结论。

% MMW-LIMITATION: D4
问题三未执行规定的储能动作候选搜索，只采用逐时规则型 SOC 前向策略。因此，成本增加以及峰值、爬坡量和碳排放不变，仅描述该已实现规则型候选。

\subsection{问题四：算—储—电联合调度与情景复算}

\subsubsection{联合负荷与服务质量}

% MMW-ID: OBJ-Q4
联合模型拟同时评价运行成本、碳排放、有效网络时延、有效等待损失、新能源未利用率及六区域非负峰值净购电功率。各维度只在有定义、数值有限且候选之间存在正尺度时进入代理评分。

% MMW-ID: CON-Q4-1
联合设施负荷由真实任务方案重新构造：
\begin{equation}
L_{rt}^{F,avg}=
\pi_r\left[
L_{rt}^{N}
+\frac{1}{\Delta t}
\sum_{i,s}g_ip_{k_i}^{GPU}\omega_{it}(s)x_{irs}
\right].
\label{eq:q4_facility_load}
\end{equation}
任务侧执行问题一全部硬约束，能源侧执行问题三的新能源守恒、负荷平衡、SOC、功率和互斥约束，外送上限取 $\min\{X_r^{grid},X_r^{storage}\}$。

% MMW-ID: CON-Q4-2
联合初始方案必须从 $S_{r,-1}=S_r^0$ 独立前推至第 $2406$ 小时，不因问题三无储能比较状态的可行性而停止。若联合构造在声明预算内没有形成完整可行候选，才返回 \texttt{feasibility\_unresolved}。

% MMW-ID: CON-Q4-3
实时即时开工率、SLA 满足率、按期完成率和等待质量均先执行分母检查。前三项被作为硬约束合规指标，代理评分只使用分母为正且候选之间存在差异的等待损失。实际程序以
\begin{equation}
Q_{\mathrm{service}}=\frac{1}{1+\overline W}
\end{equation}
作为平均等待小时 $\overline W$ 的有限代理，所得服务质量等待分量为 $1.0$。

% MMW-LIMITATION: D7
服务质量没有采用 formulation 中按最大允许等待时间归一化的完整 $Q_{\mathrm{wait}}$，而使用 $1/(1+\text{平均等待小时})$。因此，该指标只反映平均等待代理，不应被解释为完整服务质量综合评价。

\subsubsection{情景定义与实际实现}

% MMW-ID: CON-Q4-4
碳约束水平、峰谷价差系数、售电价格系数和新能源水平系数按方法契约给定的离散集合构造；新能源波动情景采用固定种子 $2026$，并通过非负截断保证新能源出力不小于零。由于截断会改变扰动分布，截断后不再假定扰动仍保持原均值或方差。

% MMW-ID: CON-Q4-5
formulation 要求联合基准与全部情景共享候选评价上限，各情景从同一联合基线独立重构，达到共享上限后未执行的情景应标记为 \texttt{not\_run\_budget\_exhausted}。实际运行共检查 $15$ 个情景候选。

% MMW-ID: CON-Q4-6
所有候选在启动前和可中断循环内部均应检查单调时钟，达到搜索截止后立即丢弃尚未完成的候选，后续时间只用于文件输出、完整审计和状态整理。实际正式任务构造使用内部截止并预留后续复算与审计时间。

% MMW-LIMITATION: D5
问题四未执行共享交替邻域搜索，每个情景只重建并审计一个规则型候选。因此，情景表反映固定规则在不同输入下的复算结果，而不是邻域搜索后的候选集合。

% MMW-LIMITATION: D6
碳约束情景没有执行额外的碳定向搜索。若基准规则型候选未达到目标，只能报告 \texttt{target\_infeasible\_within\_declared\_heuristic}，不能据此判断该碳目标在数学上不可行。

% MMW-LIMITATION: D9
只有新能源波动情景使用固定种子 $2026$ 产生随机扰动，其余任务、储能和价格情景均为确定性复算。

\subsubsection{联合基准结果}

联合基准运行成本为 $-419353923.7785073$ 元，碳排放为 $14.391601534493985$ tCO$_2$，有效网络时延为 $5.159384110987204$ ms，新能源利用率为 $0.6791293391777464$，六区域非负峰值净购电功率之和为 $40.46636711666906$ MW。联合任务方案包含 $209$ 个跨来源区域执行的任务，储能策略累计充电量为 $2076.31548844658$ MWh。任务、GPU、IT 功率、设施功率、时延、截止时间、SOC、能源平衡和购售电边界的统一审计均通过，故该方案被标记为 \texttt{full\_horizon\_feasible=true}。

需要强调的是，运行结果同时标记为 \texttt{non\_Pareto=true} 和 \texttt{global\_optimality\_certificate=false}。因此，联合方案只是在实际规则构造和已检查情景集合内通过完整约束审计的可行候选，不代表 Pareto 前沿，也不代表连续或离散决策空间中的最优方案。

\subsection{统一约束审计}

任务合法性按
\begin{align}
a_i&\le s_i,\qquad c_i\le F_i,\\
\ell_{o_ir_i}&\le\ell_i^{\max},\\
i\in\mathcal I_{\mathrm{RT}}&\Rightarrow s_i=a_i,\\
\omega_{i,2406}(s_i)&=0
\end{align}
逐项检查。问题二的能源残差为
\begin{equation}
e_{rt}^{q2}=G_{rt}^{buy}+R_{rt}^{av}
-L_{rt}^{F,avg}-G_{rt}^{sell}-R_{rt}^{curt},
\end{equation}
问题三、四的能源与 SOC 残差为
\begin{align}
e_{rt}^{energy}
&=G_{rt}^{buy}+R_{rt}^{av}+D_{rt}
-L_{rt}^{F,avg}-C_{rt}-G_{rt}^{sell}-R_{rt}^{curt},\\
e_{rt}^{SOC}
&=S_{rt}-S_{r,t-1}
-\eta_r^cC_{rt}\Delta t
+\frac{D_{rt}\Delta t}{\eta_r^d}.
\end{align}
残差按照对应硬上限或预先固定的有效输入尺度归一化。只有任务唯一执行、实时开工、截止、时延、GPU、IT、设施功率、能源守恒、SOC、终端 SOC、购售电互斥、充放电互斥、外送上限及第 $2406$ 小时禁占用全部通过时，才输出全时域可行状态。本次结果的约束满足状态为真，且 \texttt{full\_horizon\_feasible=true}。

% MMW-LIMITATION: L1
实际方法未提供最优性证书，所有结论均限定于已构造并通过审计的启发式候选。

% MMW-LIMITATION: L2
任务候选和储能动作的完整局部邻域未被穷举，故未检查的邻域中可能存在更优可行方案。

% MMW-LIMITATION: L3
碳目标未达到仅说明声明的规则型启发式候选未达到目标，不构成数学不可行性的证明。

% MMW-LIMITATION: L4
问题四各情景保持联合基准的任务迁移方案不变，情景差异主要来自能源流和储能轨迹复算，不能反映任务迁移方向对情景变化的全部响应。

% MMW-LIMITATION: L5
储能策略未计入循环退化成本，频繁充放电方案在实际工程中的经济性可能因此被高估。

% MMW-LIMITATION: L6
复杂实例可能在内部任务构造截止前返回 \texttt{feasibility\_unresolved}；该状态表示运行预算内未完成构造，而不是约束系统无解。

% MMW-LIMITATION: L7
问题二可读取区域级最大购电与最大外送功率，但不会读取储能售电限制或任何 SOC 字段，该边界差异应在跨问题比较时保留。

% MMW-LIMITATION: L8
试跑文件只证明真实输入基础门禁和有限样本时间窗检查通过，不能替代正式全时域方法的可行性审计。
